\documentclass{article}
\usepackage[UTF8]{ctex}
\usepackage{amsmath,amssymb,amsfonts}
\usepackage{enumitem}
\usepackage{geometry}
\geometry{a4paper, margin=1in}
% === 错题本排版优化（不参与题目标识） ===
\setlength{\parindent}{0pt}
\setlength{\parskip}{0.35em plus 0.10em minus 0.08em}
\setlength{\emergencystretch}{3em}
\setlist{itemsep=0.25em, topsep=0.40em, parsep=0pt}
\allowdisplaybreaks[2]
% === 错题本排版优化结束 ===
\begin{document}
\section*{湘潭大学2026年数学分析真题}
\section*{198．湘潭大学2026年数学分析试卷}

\subsection{题目1}
1．（30 分）计算题．

\begin{itemize}
  \item[（1）] （10 分）求极限 \(\lim _{x \rightarrow 0}\left(\cos x-\frac{x^{2}}{2}\right)^{\frac{1}{x^{2}}}\) ．
  \item[（2）] （10 分）求极限 \(\lim _{x \rightarrow 0} x\left[\frac{1}{x}\right]\) ．
  \item[（3）] （10 分）求积分 \(\int_{0}^{\frac{\pi}{2}} \frac{\sin ^{2} x}{\sin x+\cos x} \mathrm{~d} x\) ．
\end{itemize}

\subsection{题目2}
2．（20 分）计算题．

\begin{itemize}
  \item[（1）] （10 分）设 \(y=\sqrt[3]{x-\tan x}\) ，求 \(\mathrm{d} y\) 及 \(\mathrm{d}^{2} y\) ．
  \item[（2）] （10 分）设 \(F(y)=\int_{a}^{b} f(x)|y-x| \mathrm{d} x\) ，其中 \(f(x)\) 可微，求 \(F^{\prime \prime}(y)\) ．
\end{itemize}

\subsection{题目3}
3．（15 分）计算重积分 \(\iiint_{V} \sqrt{1-\frac{x^{2}}{a^{2}}-\frac{y^{2}}{b^{2}}-\frac{z^{2}}{c^{2}}} \mathrm{~d} x \mathrm{~d} y \mathrm{~d} z\) ，其中 \(V=\left\{(x, y, z) \left\lvert\, \frac{x^{2}}{a^{2}}+\frac{y^{2}}{b^{2}}+\frac{z^{2}}{c^{2}} \leq 1\right.\right\}\) ．

\subsection{题目4}
4．（15 分）计算曲面积分 \(\iint_{\sum} \frac{x \mathrm{~d} y \mathrm{~d} z+y \mathrm{~d} z \mathrm{~d} x+z \mathrm{~d} x \mathrm{~d} y}{\left(x^{2}+y^{2}+z^{2}\right)^{\frac{3}{2}}}\) ，其中 \(\sum: x^{2}+2 y^{2}+3 z^{2}=1\) ，取外侧．

\subsection{题目5}
5．（20 分）讨论广义积分 \(I=\int_{0}^{1} x^{p-1} \ln ^{2} x \mathrm{~d} x\) 的一致收敛性：

\begin{itemize}
  \item[（1）] （10 分）\(p \geq p_{0}>0\) ．
  \item[（2）] （10 分）\(p>0\) ．
\end{itemize}

\subsection{题目6}
6．（20 分）设 \(f(x)\) 在 \([a, b]\) 上连续．

\begin{itemize}
  \item[（1）] （10 分）证明：\(f(x)\) 在 \([a, b]\) 上有界．
  \item[（2）] （10 分）证明：\(f(x)\) 在 \([a, b]\) 上一定有最大值和最小值．
\end{itemize}

\subsection{题目7}
7．（20 分）设 \(S_{n}(x)=n^{\alpha} x e^{-n x}\) ，其中 \(\alpha\) 是参数，求 \(\alpha\) 的取值范围，使得函数列 \(\left\{S_{n}(x)\right\}\) 在［0，1］上：

\begin{itemize}
  \item[（1）] （6 分）一致收敛。
  \item[（2）] （6 分）积分运算与极限运算可以交换，即 \(\lim _{n \rightarrow \infty} \int_{0}^{1} S_{n}(x) \mathrm{d} x=\int_{0}^{1} \lim _{n \rightarrow \infty} S_{n}(x) \mathrm{d} x\) ．
  \item[（3）] （8 分）求导运算与极限运算可以交换，即对一切 \(x \in[0,1]\) ，成立 \(\lim _{n \rightarrow \infty} \frac{\mathrm{~d}}{\mathrm{~d} x} S_{n}(x)=\frac{\mathrm{d}}{\mathrm{d} x} \lim _{n \rightarrow \infty} S_{n}(x)\) ．
\end{itemize}

\subsection{题目8}
8．（10 分）求 \(\cos x\) 的全部零点的倒数平方之和。
\end{document}
